\documentclass[11pt]{amsart}
\usepackage[T1]{fontenc}
\usepackage[utf8]{inputenc}
\usepackage[french]{babel}
\usepackage{amssymb}
\usepackage[margin=1.05in]{geometry}
\usepackage{xurl}
\usepackage[hidelinks]{hyperref}
\newtheorem{lemma}{Lemme}
\newtheorem{proposition}[lemma]{Proposition}
\newcommand{\OO}{\mathcal O}
\newcommand{\m}{\mathfrak m}
\title{Du relèvement fini au cas étale universel à deux variables}
\date{26 octobre 2026; complément du 1 novembre 2026}
\begin{document}
\maketitle

\noindent\textbf{Statut et apport.}
Cette note explique comment remplacer la fibre particulière
$p=x$, $q=y$ par une fibre étale finie quelconque dans l'argument
de \texttt{finite-lifting.tex}. Le point nouveau est la construction,
sur les trois cartes de $\mathbf P^2$, des sections approchées requises
par ce lemme. Elle utilise seulement des calculs dans des algèbres
finies sur le corps résiduel et le relèvement polynomial des idempotents.
La proposition ci-dessous est conditionnelle à deux données de finitude
cohomologique explicitement énoncées. On n'affirme pas que leur
construction universelle sur $\mathbf Z$, ni le présent raccord,
soient déjà formalisés dans Unit. Le complément final raccorde cette
construction à la fin locale raccourcie désormais vérifiée dans Unit.
Il montre aussi que la seconde donnée, la finitude de tout
$H^0(W,\OO_W)$, peut être supprimée : un module fini construit à partir
de la section corrigée suffit. Ce raccord complet reste à formaliser.

\section{Une localisation finie sur un corps est un facteur direct}

\begin{lemma}\label{factor}
Soient $k$ un corps discret, $C$ une $k$-algèbre et $f\in C$.
Supposons donnée une structure d'algèbre de dimension finie sur
$L=C[1/f]$, avec les données d'algèbre linéaire nécessaires.
On peut construire $e\in C$ tel que
\[
 e^2=e,\qquad C[1/e]\simeq C[1/f].
\]
L'algèbre $Ce$, d'unité $e$, s'identifie à $L$.
\end{lemma}
\begin{proof}
Si $L=0$, prendre $e=0$. Sinon, la multiplication par $f$ dans $L$
est une matrice inversible. Son polynôme caractéristique a un terme
constant non nul; Cayley--Hamilton exprime donc $f^{-1}$ comme un
polynôme en $f$ à coefficients dans $k$. Soit $b\in k[f]\subseteq C$
un tel représentant. L'égalité $1-fb=0$ dans $C[1/f]$ fournit un
entier $N\geq1$ et le certificat
\[
 f^N(1-fb)=0\quad\hbox{dans }C.
\]
Poser $e=(fb)^N$. Alors $e(1-fb)=0$, donc $e^2=e$.
Dans $Ce$, l'élément $fe$ est inversible, d'inverse $be$.
Dans $C[1/f]$, l'image de $e$ est $1$.
Les deux propriétés universelles donnent les isomorphismes annoncés.
La construction ne suppose pas $C$ de dimension finie.
\end{proof}

Une localisation d'une algèbre de dimension finie sur $k$ est encore
de dimension finie. Cela se calcule par algèbre linéaire, en séparant
le noyau stabilisé et l'image stabilisée de la multiplication par
l'élément inversé. Il n'est pas nécessaire de factoriser des polynômes
en facteurs irréductibles.

\section{Les données résiduelles sur les trois cartes}

Soient $(R,\m,k)$ un anneau local résiduellement discret,
$B=R[x,y]/(F,G)$ et $E=B[1/\Delta]$, où $\Delta\in R[x,y]$.
On suppose que $\bar E=E/\m E$ est une $k$-algèbre de dimension finie
effectivement présentée. Pour l'application étale, $\Delta$ est le
déterminant jacobien; la finitude de $\bar E$ est le résultat
constructif sur les algèbres de type fini non ramifiées sur un corps.
Ce résultat sur un corps est une entrée distincte, et non une
conséquence du ZMT que l'on cherche à remplacer.

Homogénéiser $F,G$ en $P,Q$ et poser
\[
 W=\operatorname{Proj}R[X,Y,Z]/(P,Q).
\]
Choisir un homogénéisé $\Delta^h$ de degré $d$ de $\Delta$.
Pour $L=X,Y,Z$, soit $C_L$ l'anneau de la carte $D_+(L)$ de $W$,
et poser
\[
 f_L=\frac{Z\Delta^h}{L^{d+1}}\in C_L.
\]
Sur la fibre résiduelle, on a
\[
 (\bar C_X)_{\bar f_X}\simeq\bar E[1/x],\quad
 (\bar C_Y)_{\bar f_Y}\simeq\bar E[1/y],\quad
 (\bar C_Z)_{\bar f_Z}\simeq\bar E.
\]
Ces trois algèbres sont de dimension finie sur $k$. Le lemme
\ref{factor} donne donc des idempotents $\bar e_L\in\bar C_L$.
Ils se recollent sur les intersections : ils définissent le même
ouvert $D_+(Z\Delta^h)$, et l'idempotent correspondant à un facteur
ouvert et fermé est unique. Cette unicité se vérifie algébriquement
par la décomposition en deux facteurs associée à un idempotent.

Dans le facteur $\bar e_L\bar C_L$, $Z/L$ est inversible : le
produit $(Z/L)(\Delta^h/L^d)$ y est inversible.
Construire son inverse donne une section
\[
 \bar u_L\in\Gamma(D_+(L),\OO_{W_k}(-1)),\qquad
 Z\bar u_L=\bar e_L,\quad \bar e_L\bar u_L=\bar u_L.
\]
Sur la carte $Z$, c'est simplement $\bar e_Z/Z$.
Les $\bar u_L$ se recollent par unicité de cet inverse dans chaque
facteur. Leur restriction à $\operatorname{Spec}\bar E$ est $1/Z$.
On a donc construit une section résiduelle de $\OO(-1)$ supportée
sur la partie affine finie, sans choisir de points géométriques,
de racines, ou de décomposition en corps.

\section{Un idéal fini suffit}

Choisir des représentants $e_L,u_L$ sur les trois cartes de $W$.
Les défauts des identités
\[
 e_L^2=e_L,\quad Zu_L=e_L,\quad e_Lu_L=u_L,
\]
et les défauts de recollement appartiennent à $\m$ fois les modules
affines correspondants. Chaque appartenance possède un témoin fini.
Rassembler les coefficients de $\m$ figurant dans ces témoins donne
un idéal \emph{fini} $I=(a_1,\ldots,a_r)\subseteq\m$ tel que toutes
ces identités soient déjà vraies modulo $I$.
Cette étape évite de supposer que l'on sait calculer des générateurs
de l'idéal maximal $\m$.

\section{Relèvement explicite à chaque ordre}

\begin{lemma}\label{nilpotent}
Pour chaque $n\geq1$, les données précédentes fournissent une section de
\[
 \OO_W(-1)/I^n\OO_W(-1)
\]
qui relève la section modulo $I$.
\end{lemma}
\begin{proof}
Travailler sur chaque carte modulo $I^n$.
Partant de $e_L$, itérer
\[
 t\longmapsto 3t^2-2t^3.
\]
Si $\delta=t^2-t$, le nouveau défaut est
$-\delta^2(3-2t)(2t+1)$. Après $r$ itérations avec $2^r\geq n$,
on obtient un idempotent $e_{L,n}$ qui relève $e_L\bmod I$.
La formule fonctionne en toute caractéristique.

Poser $v=e_{L,n}u_L$ et $w=e_{L,n}-Zv$. Alors $w\in I C_L/I^nC_L$,
donc $w^n=0$, et
\[
 u_{L,n}=v(1+w+\cdots+w^{n-1})
\]
satisfait $Zu_{L,n}=e_{L,n}$ et $e_{L,n}u_{L,n}=u_{L,n}$.
Les idempotents qui relèvent un même idempotent modulo un idéal
nilpotent sont uniques dans un anneau commutatif : si $e,e'$ sont
deux tels relèvements, $e(1-e')$ et $e'(1-e)$ sont des idempotents
nilpotents, donc nuls. Ils se recollent donc sur les intersections.
Les inverses de $Z$ sur leurs facteurs sont également uniques;
les $u_{L,n}$ se recollent.
\end{proof}

Pour le complexe de Čech $K^\bullet$ de $\OO_W(-1)$, choisir les
représentants polynomiaux de ces formules donne des $z_n\in K^0$ avec
\[
 dz_n\in I^nK^1,\qquad z_n-z_1\in IK^0.
\]
Ce sont exactement les données du lemme de correction finie.
Contrairement au cas $p=x$, $q=y\bmod I$, on ne prétend pas utiliser
la même cochaîne à tous les ordres.

\section{Le résultat conditionnel précis}

\begin{proposition}\label{hull}
Dans la situation précédente, supposons fournies les deux données :
\begin{enumerate}
\item le module gradué
$M=\bigoplus_{n\geq0}H^1(I^nK^\bullet)$ possède la propriété
noethérienne constructive utilisée dans \texttt{finite-lifting.tex};
\item $H^0(W,\OO_W)$ possède une famille finie de générateurs comme
$R$-module, avec sa propriété de génération.
\end{enumerate}
On peut construire une sous-algèbre $D\subseteq E$, finie sur $R$,
et $t\in D\cap(1+\m E)$ tels que $D[1/t]=E[1/t]$.
\end{proposition}
\begin{proof}
La correction finie appliquée aux $z_n$ donne
$h\in H^0(W,\OO_W(-1))$ de la réduction prescrite.
Les trois sections $Zh,Xh,Yh$ appartiennent à l'algèbre finie
$H^0(W,\OO_W)$. Leurs images dans $B$, puis dans $E$, sont
$s,sx,sy$, et $s\equiv1\bmod\m E$.
Le déterminant de multiplication sur une famille génératrice donne
une équation monique pour chacune. Ainsi
\[
 D=R[s,sx,sy]\subseteq E
\]
est fini sur $R$. Aucune injectivité de la restriction
$H^0(W,\OO_W)\to B$ n'est requise.

Dans $\bar E$, l'inverse de $\Delta$ est un polynôme en $\Delta$,
par le même calcul de Cayley--Hamilton que précédemment. Choisir
un relèvement $b\in R[x,y]$ de ce polynôme; alors
$\Delta b\equiv1\bmod\m E$.
Pour $N\geq\deg(\Delta b)+1$, poser
\[
 t=s^N\Delta b.
\]
On a $t\in D$ et même $t=s\,(s^{N-1}\Delta b)$ avec le second
facteur dans $D$. Donc $s$ est inversible dans $D[1/t]$, qui contient
$x=(sx)/s$, $y=(sy)/s$ et
$\Delta^{-1}=s^Nb/t$.
Il contient par conséquent tous les générateurs de $E[1/t]$.
Enfin $t\equiv1\bmod\m E$, comme requis.
\end{proof}

\section{Ce que cela change pour le cas universel}

Fixer des bornes de degrés et prendre
\[
 R_0=\mathbf Z[c_\alpha,d_\beta],\quad
 F=\sum c_\alpha X^{\alpha_1}Y^{\alpha_2},\quad
 G=\sum d_\beta X^{\beta_1}Y^{\beta_2},\quad
 E_0=(R_0[X,Y]/(F,G))[1/\Delta].
\]
On doit travailler sur $R_0$, et non seulement sur son corps des
fractions : une preuve sur ce dernier perdrait les spécialisations
où les dénominateurs s'annulent.

Après localisation de $R_0$ à la contraction d'un premier décidable
générique de $E_0$, la construction ci-dessus donne exactement la
conclusion du corollaire 1.4 d'ACL utilisée dans
\texttt{zmt2-revised.tex} (dans le répertoire \texttt{COLL}).
Elle peut donc remplacer \emph{cet appel} au ZMT, dès que la donnée
de correction finie est disponible : la proposition~\ref{hull-short}
ci-dessous supprime la seconde hypothèse de finitude.
Il reste à rendre l'algèbre monogène localement par un déterminant;
la fin locale raccourcie fournit alors la relation monique simple et
la suppression du noyau. Pour obtenir un recouvrement principal fini
du cas universel, il faut encore assembler le passage local--global
et la descente des certificats finis. Un tel recouvrement se
spécialiserait à tout anneau de coefficients.

Les calculs de \texttt{finite-cohomology.tex} indiquent comment fournir
les deux données de finitude : présentation par élimination du module
de Rees pour $H^1$, puis noyau d'une matrice finie pour $H^0$ dans le
cas de deux homogénéisés. Le calcul monomial de Čech est valable
sur tout anneau. En revanche, la construction effectivement vérifiée
des présentations et noyaux sur $\mathbf Z[c_\alpha,d_\beta]$, et leur
transport aux localisations, ne découle pas de la seule déclaration
de noethérianité déjà disponible dans Unit. La référence
Coquand--Persson sur un corps ne couvre pas à elle seule ce passage
sur $\mathbf Z$.

\medskip
\noindent\textbf{Travail de formalisation désormais délimité.}
Les nouveaux lemmes à isoler sont : localisation finie sur un corps
$\Rightarrow$ idempotent avec certificat; calcul des trois cartes
résiduelles; extraction de l'idéal fini $I$; relèvement polynomial et
recollement des inverses; raccord au lemme de correction; enfin le
changement de dénominateur $t=s^N\Delta b$.
La proposition~\ref{hull} est une preuve sur papier sous les deux
entrées explicites de finitude, pas encore une preuve Unit du théorème
universel. Elle lève toutefois l'obstacle particulier de la fibre
résiduelle non rationnelle et de l'inversion du jacobien, sans ajouter
une variable à la compactification.

\section{Raccord avec la fin locale raccourcie}

Le relèvement de cette note et le nouvel argument de
\path{etale-monic-cover-analysis.tex} interviennent à deux endroits
différents de la preuve locale. Le premier construit une sous-algèbre
finie; le second transforme des données de génération monogène
intégrale en une présentation standard monique. Leur combinaison
permet aussi de simplifier l'étape intermédiaire de finitude.

\subsection{La finitude de tout le module des sections est superflue}

\begin{proposition}\label{hull-short}
Dans la proposition~\ref{hull}, la première donnée de finitude suffit.
Il n'est pas nécessaire de supposer que $H^0(W,\OO_W)$ est de type
fini sur $R$.
\end{proposition}
\begin{proof}
La correction fournit toujours
$h\in H^0(W,\OO_W(-1))$ avec la réduction prescrite. Sur la carte
$Z$, écrire $h=s/Z$, avec $s\in B$. Choisir un représentant
polynomial $S(x,y)$ de $s$ de degré au plus $d_0$.
Sur la carte $X$, un représentant de $h$ s'écrit
$A(X,Y,Z)/X^{a+1}$, avec $A$ homogène de degré $a$.
L'égalité sur l'intersection avec la carte $Z$ donne dans $B[1/x]$
\[
                  s x^{a+1}=A(x,y,1).
\]
Une puissance $x^r$ annule le défaut de cette égalité dans $B$.
Après multiplication par cette puissance, on obtient donc dans $B$
une relation
\[
               s x^{a_x+1}=q_x(x,y),\qquad \deg q_x\leq a_x,
               \qquad a_x=a+r.
\]
La carte $Y$ fournit de même
\[
               s y^{a_y+1}=q_y(x,y),\qquad \deg q_y\leq a_y.
\]
Ces égalités ne supposent pas que $B$ s'injecte dans ses localisations.
Les puissances qui annulent les défauts font partie des témoins.

Poser $d=d_0+a_x+a_y$ et prendre pour $V\subseteq B$ le module
engendré par les monômes $x^iy^j$ avec $i+j\leq d$.
Ce module contient $1$ et est stable par $s,sx,sy$.
En effet, pour réduire $s x^iy^j$ avec $i+j\leq d+1$, deux cas
finis suffisent. Si $i+j\leq a_x+a_y$, utiliser $S$ : le degré
obtenu est au plus $d$. Sinon $i\geq a_x+1$ ou $j\geq a_y+1$;
la relation correspondante diminue le degré total d'au moins un,
qui devient donc au plus $d$.

Le déterminant de multiplication sur cette famille finie, qui contient
$1$, donne des équations moniques pour $s,sx,sy$. Une base ou une
indépendance linéaire n'est pas nécessaire. Leur sous-algèbre image
\[
                         D=R[s,sx,sy]\subseteq E
\]
est donc finie sur $R$. La réduction prescrite donne
$s\equiv1\pmod{\m E}$, car $\bar e_Z$ devient $1$ dans $\bar E$.
Le changement de dénominateur $t=s^N\Delta b$ de la preuve précédente
donne alors $D[1/t]=E[1/t]$ et $t\in1+\m E$, sans aucun recours
à la finitude de $H^0(W,\OO_W)$.
\end{proof}

Cette réduction correspond à la construction déjà vérifiée dans
\path{sag/zmt-certificate-extract-general}. Le rôle de la section
fermée est de fournir des relations qui abaissent le degré. La finitude
requise porte seulement sur le module explicite $V$ construit à partir
de ces relations, et non sur toutes les sections de degré zéro.

\subsection{Le raccord exact après le déterminant}

Supposons maintenant $E$ formellement étale et fixons le voisinage
résiduel auquel on veut appliquer l'argument local. Les données
monogènes de ce facteur résiduel et le déterminant de la preuve locale
donnent un élément $\eta\in D$ et
\[
                p\in R[\eta],\qquad pD\subseteq R[\eta].
\]
Le choix résiduel garantit que $p$ n'appartient pas au premier choisi;
$t$ n'y appartient pas non plus puisque $t\equiv1\pmod{\m E}$.
Comme $t\in D$, on peut choisir $\Gamma\in R[T]$ avec
\[
                            \Gamma(\eta)=pt.
\]
L'égalité $D[1/t]=E[1/t]$ fournit, pour chaque $e\in E$, un entier
$n$ et $b\in D$ avec $t^ne=b$, après avoir multiplié par une
puissance supplémentaire de $t$ si nécessaire. Alors
\[
                 \Gamma(\eta)^{n+1}e=p^n\,p(tb)\in R[\eta].
\]
C'est exactement l'hypothèse de génération par fractions du nouveau
théorème local. Comme $D$ est fini, $\eta$ possède aussi une relation
monique $Q_0(\eta)=0$, obtenue par le déterminant de multiplication.
Il n'est pas nécessaire que $D$ soit de présentation finie ou libre
sur $R$.

On applique maintenant la fin raccourcie à
$E[1/\Gamma(\eta)]$. La présentation finie et la non-ramification
formelle construisent un polynôme $H$ et un entier $N$ tels que
\[
              H(\eta)=0,\qquad H'(\eta)=\Gamma(\eta)^N.
\]
Si $k$ est la longueur d'une liste de coefficients de $H$, le polynôme
\[
                          Q=H+T^kQ_0^2
\]
est monique, annule $\eta$ et a une dérivée inversible dans cette
localisation. La suppression du noyau idempotent donne une présentation
monique standard du même $E[1/\Gamma(\eta)]$.
La formalité étale vient directement du jacobien inversible dans le
cas de cette note; pour une algèbre de présentation finie plate et
formellement non ramifiée, elle est fournie par le théorème direct
vérifié dans \path{sag/etale-direct}.

La suite combinée est donc
\[
\begin{gathered}
\text{sections approchées et correction finie}\\
\Downarrow\\
\text{relations abaissant le degré, puis }D\text{ fini}\\
\Downarrow\\
\text{générateur résiduel et déterminant : }\eta,\Gamma,Q_0\\
\Downarrow\\
\text{relation simple monique et suppression du noyau.}
\end{gathered}
\]
Elle supprime la donnée globale de finitude de $H^0$, le calcul spécial
destiné à comparer une relation résiduelle à $Q_0$, la boucle
gcd/Nakayama et le rétrécissement supplémentaire par la dérivée sur
l'algèbre cible. Le choix d'un générateur d'un facteur résiduel reste
nécessaire pour le déterminant; la finitude seule d'une fibre étale
ne fournit pas un générateur global, comme le montre
$\mathbf F_2^3$.

\subsection{État exact du raccord dans Unit}

Les éléments vérifiés et le travail de raccord se distinguent ainsi.
\begin{itemize}
\item La fin locale, à partir de $\eta,\Gamma$, de l'intégralité et de
la génération par fractions, est vérifiée dans
\path{sag/etale-monic-local/ShortLocalEtale.seg}. Elle n'a aucune
hypothèse de corps, d'anneau local ou de premier.
\item L'extraction des relations de degré et le module stable fini
sont vérifiés dans les développements d'extraction de ZMT. L'adaptation
à la section de la présente note reste à écrire : sur la carte $Z$,
sa réduction est $\bar e_Z/Z$, et non nécessairement $1/Z$ sur toute
la carte. L'idempotent devient $1$ seulement après restriction à
$\bar E$. Le théorème actuel
\path{GenericClearedExtraction} exige une normalisation exacte par
$1/Z$. Il faut séparer l'extraction des relations, qui utilise la
fermeture, de la preuve $s-1\in\m E$, qui utilise ici cette restriction.
\item Les trois cartes résiduelles, leurs idempotents, les relèvements
polynomiaux et leur raccord à la correction doivent encore être
assemblés pour cette compactification dans $\mathbf P^2$. La première
hypothèse cohomologique de la proposition~\ref{hull} reste explicite
dans cette voie. Les théorèmes de dépendance pondérée et de comparaison
de Rees déjà vérifiés fournissent les outils; ils ne constituent pas
automatiquement cette nouvelle instanciation.
\item Le ZMT algébrique en toute dimension est désormais vérifié dans
\path{sag/zmt-general}. Pour terminer au plus vite l'intégration locale,
on peut obtenir $D,t$ par ce théorème et conserver exactement la même
suite après le déterminant. La construction de la présente note est
une autre voie vers $D,t$, qui garde la compactification à trois cartes
malgré l'inversion de $\Delta$.
\end{itemize}

Il s'agit toujours d'un argument local. Le passage à un recouvrement
principal fini sur la base universelle, la descente des certificats et
leur spécialisation sont des tâches distinctes. En particulier
$t\in1+\m E$ n'implique pas que $t$ soit une unité de $E$; la fibre
fermée peut manquer des composantes de $E$. Aucune conclusion globale
de recouvrement n'est ajoutée par le présent raccord.

\section{Deux variables : un premier cas borné et l'obstacle restant}

Pour analyser le recouvrement, la cible précise est la suivante : pour
\[
 E=\bigl(R[X,Y]/(F,G)\bigr)[1/\Delta],\qquad
 \Delta=F_XG_Y-F_YG_X,
\]
construire une famille finie $h_\nu,c_\nu\in E$ avec
$\sum_\nu c_\nu h_\nu=1$, et une présentation monique standard de
chaque $E[1/h_\nu]$ sur $R$. Le cas où l'une des équations est linéaire
donne déjà une construction bornée sans données résiduelles.

\begin{proposition}\label{linear-two}
Supposons $F=aX+bY+c$, avec $a,b,c\in R$, et que le degré total
de $G$ soit au plus $d$, où $d\geq2$. Alors $E$ admet un recouvrement
par au plus $4d^2$ algèbres moniques standard étales sur $R$.
Les polynômes moniques peuvent tous avoir le degré $d$.
Les cartes nulles peuvent être conservées.
\end{proposition}
\begin{proof}
La formule $\Delta=aG_Y-bG_X$ donne déjà le certificat
\[
                     1=a\frac{G_Y}{\Delta}
                         +b\frac{-G_X}{\Delta}
                     \quad\hbox{dans }E.
\]
Il suffit donc de traiter les ouverts $D(a)$ et $D(b)$.
Écrire $G=\sum_{i+j\leq d}g_{ij}X^iY^j$ et définir, sur la base
originale $R$, les deux polynômes
\begin{align*}
 g_a(U)&=\sum_{i+j\leq d}g_{ij}a^{d-i}(-bU-c)^iU^j,\\
 g_b(U)&=\sum_{i+j\leq d}g_{ij}b^{d-j}U^i(-aU-c)^j.
\end{align*}
Ils ont degré au plus $d$. Les exposants sont tous non négatifs;
aucun coefficient n'a été divisé. Dans $R[X,Y]/(F)$, on a
\begin{align*}
 g_a(Y)&=a^dG(X,Y),& g'_a(Y)&=a^{d-1}\Delta(X,Y),\\
 g_b(X)&=b^dG(X,Y),& g'_b(X)&=-b^{d-1}\Delta(X,Y).
\end{align*}
Ces identités se prouvent terme à terme à partir de
$aX=-bY-c$ et $bY=-aX-c$. Elles sont vraies avant toute
localisation, même si $a$ ou $b$ est un diviseur de zéro.

Poser
\[
 C_a=(R[U]/(g_a))[1/g'_a],\qquad
 C_b=(R[U]/(g_b))[1/g'_b].
\]
Les substitutions $Y=U$, $X=-(bU+c)/a$, et respectivement
$X=U$, $Y=-(aU+c)/b$, donnent des isomorphismes d'algèbres
\[
                         E[1/a]\simeq C_a[1/a],\qquad
                         E[1/b]\simeq C_b[1/b].
\]
L'inversibilité de la dérivée correspond exactement à celle du
jacobien sur l'ouvert considéré, par les identités précédentes.

Le théorème à une équation de
\path{etale-monic-cover-analysis.tex} fournit pour $C_a$ au plus
$2d^2$ cartes moniques de degré $d$, de dénominateurs polynomiaux
$h_{a\alpha}(U)$; il en fournit autant pour $C_b$.
Une localisation supplémentaire au coefficient constant $a$, ou $b$,
conserve la présentation standard et son polynôme monique.
Les cartes correspondantes de $E$ sont donc
\[
          E[1/(a h_{a\alpha}(y))],\qquad
          E[1/(b h_{b\beta}(x))].
\]

Vérifions le recouvrement par des identités finies sur $E$.
En chassant les dénominateurs dans le certificat de recouvrement de
$C_a$, on obtient un entier $N\geq0$ et des polynômes $A_\alpha$
tels que
\[
               (g'_a)^N=\sum_\alpha A_\alpha h_{a\alpha}
                         \pmod{g_a}.
\]
En évaluant dans $E$, puis en utilisant l'inverse de $\Delta$ et
en multipliant par $a$, il vient
\[
                  a^{(d-1)N+1}\in
                       (a h_{a\alpha}(y))_\alpha.
\]
De même une puissance strictement positive de $b$ appartient à
l'idéal engendré par les $b h_{b\beta}(x)$.
Si ces deux puissances sont $a^P$ et $b^Q$, développer la puissance
$P+Q-1$ du certificat $1=aG_Y/\Delta-bG_X/\Delta$ : chaque
terme contient $a^P$ ou $b^Q$. La substitution des deux témoins
fournit l'identité de recouvrement cherchée. C'est un calcul fini,
sans recours à des premiers ou à une décision d'inversibilité.
\end{proof}

La proposition est ici un résultat mathématique, conditionné seulement
au théorème à une équation démontré dans l'autre note; leur
formalisation n'est pas revendiquée. Elle fournit un premier test
de l'ensemble du passage à un recouvrement, y compris la chasse aux
dénominateurs et le maintien de la base $R$.

\subsection{La difficulté du recouvrement existe déjà à deux variables}

Même deux équations quadratiques peuvent cacher une composante à la
fibre fermée. Soient $R=k[t]_{(t)}$ et
\[
 F=X^2-X,\qquad G=(1-X)Y+X(tY-1).
\]
Dans $B=R[X,Y]/(F,G)$, l'idempotent $x$ donne
\[
                         B\simeq R\times R[1/t].
\]
Le jacobien est
\[
              \Delta=(2x-1)(1-x+tx),
\]
et son inverse est déjà dans $B$ :
\[
              \Delta^{-1}=(2x-1)(1-x+xy).
\]
En effet $(2x-1)^2=1$ et $txy=x$, donc
$(1-x+tx)(1-x+xy)=1$. Ainsi $E=B[1/\Delta]=B$.
L'idempotent $1-x$ satisfait
\[
                       1-x\equiv1\pmod{tE},
\]
mais son ouvert ne contient que le facteur $R$. Le recouvrement
complet conserve les deux ouverts $D(1-x)$ et $D(x)$, avec le
certificat $(1-x)+x=1$; le second est standard sur $R$ avec polynôme
monique $T$ et dénominateur $t$.

Cet exemple situe exactement la limite du seul relèvement de la fibre
fermée. Réduire le nombre de variables à deux ne supprime pas cette
difficulté. Il faut garder un certificat de recouvrement, et non
seulement un élément de $1+\m E$.

\subsection{La prochaine question pour deux équations non linéaires}

Lorsque les deux équations sont non linéaires, l'identité jacobienne
fournit encore le recouvrement $D(F_X)\cup D(F_Y)$ de $E$.
Cependant l'inversibilité de $F_X$ n'exprime pas $x$ comme un
polynôme ou une fraction en $y$ sur la base originale. Elle donne
une résolution formelle à travers des idéaux nilpotents; l'élimination
algébrique utilisée dans la proposition~\ref{linear-two} reste à
construire.

L'objectif intermédiaire précis est donc une famille finie de données
monogènes intégrales, avec des dénominateurs dont on prouve qu'ils
engendrent l'idéal unité. Dès que ces données sont disponibles, le
théorème local raccourci fournit les présentations standard.
La construction projective à trois cartes fournit une voie locale
vers ces données; elle ne produit pas encore cette famille couvrante.
La proposition~\ref{linear-two} isole un cas où l'on sait produire
directement les données et le certificat de recouvrement.

À titre de vérification arithmétique, les identités d'élimination et
de dérivée ont été testées, monôme par monôme, pour les bornes de
degré $1$ à $4$ et sur $\mathbf Z/m\mathbf Z$ avec
$m\in\{2,3,4,5,6,8\}$ : $216\,580$ cas passent. Les identités du
dernier exemple passent pour les $50$ solutions de ses équations dans
ces mêmes anneaux. Ces calculs contrôlent les signes et les puissances;
ils ne remplacent ni les preuves ci-dessus ni une vérification Unit.

\end{document}
